\section{The Angels and the Order of the Universe}\label{sec:order-universe}

God made the world to show forth his goodness, and he made it many and
unequal so that it might show that goodness better. ``He brought things into
being in order that His goodness might be communicated to creatures, and be
represented by them; and because His goodness could not be adequately
represented by one creature alone, He produced many and diverse creatures,
that what was wanting to one in the representation of the divine goodness
might be supplied by another. For goodness, which in God is simple and
uniform, in creatures is manifold and divided and hence the whole universe
together participates the divine goodness more perfectly, and represents it
better than any single creature whatever'' (\ST{I q.~47 a.~1}). The same
wisdom that distinguishes things makes them unequal, ``for the universe would
not be perfect if only one grade of goodness were found in things'' (\ST{I
q.~47 a.~2}); Ecclesiasticus asks ``Why doth one day excel another, and one
light another \dots\ when all come of the sun?'' and answers, ``By the
knowledge of the Lord they were distinguished'' (Ecclus 33:7--8). And the
many unequal things make one world, ``for this world is called one by the
unity of order, whereby some things are ordered to others'' (\ST{I q.~47
a.~3}). The whole of the angelic treatise stands inside these three articles.
The angels are the highest grade of the goodness God communicates, the first
links of the order by which the many are one, and the first voices of the
praise in which the order returns to him.

\subsection{The Summit of Creation}

The \emph{Summa contra gentiles} opens its account of the distinct creatures
with the intellectual creatures and gives the reason in its title:
\latin{Quod oportuit ad perfectionem universi aliquas creaturas
intellectuales esse}. By the divine disposition, which assigns to created
things the best perfection in their mode, it followed that certain
intellectual creatures should be made, \latin{in summo rerum vertice
constitutae}, set at the very summit of things (\emph{SCG} II.46). Aquinas
gives the reasons one after another, and together they are a theology of the
angels in miniature. An effect is most perfect when it returns to its
principle, \latin{unde et circulus inter omnes figuras, et motus circularis
inter omnes motus, est maxime perfectus}; all creatures return to God by
bearing his likeness in their being and nature; but since God is the
principle of creatures by his intellect, it was necessary for their
perfection that some creatures should understand. The first perfection of a
thing is its being, the second its operation, and the consummate perfection
of the universe required creatures that return to God \latin{non solum
secundum naturae similitudinem, sed etiam per operationem}, which can only
be by intellect and will, \latin{quia nec ipse Deus aliter erga seipsum
operationem habet}, since God himself has no other operation toward
himself. A thing is perfectly like another in acting when it acts not only
with the same kind of action but in the same mode; God acts by intellect and
will; there must therefore be creatures that act as God acts. The form by
which God makes the creature is an intelligible form in his mind; the
highest likeness requires creatures in which the form of the divine
intellect is expressed \latin{secundum esse intelligibile}. And as God
contains all creatures in himself, not by extension but simply, so
\latin{factae sunt creaturae intellectuales, quae creaturas corporales
continerent, non extensione quantitatis, sed simpliciter per modum
intelligibilem} (\emph{SCG} II.46). The \emph{Summa} gathers the argument
into two sentences: ``God produces the creature by His intellect and will
\dots\ Hence the perfection of the universe requires that there should be
intellectual creatures. Now intelligence cannot be the action of a body, nor
of any corporeal faculty; for every body is limited to `here' and `now.'
Hence the perfection of the universe requires the existence of an incorporeal
creature'' (\ST{I q.~50 a.~1}; \S\ref{sec:q50}).

Dionysius sets the same summit in the language of participation. ``It was
through goodness that the superessential Godhead, having fixed all the
essences of things being, brought them into being,'' and all things share in
a Providence that ``bubbles forth'' from it: things without life by their
being, living things by its life-giving power, ``things rational and
intellectual participate in its self-perfect and preeminently perfect wisdom,
above all reason and mind'' (\emph{CH} 4.1). ``The holy orders, then, of the
Heavenly Beings share in the supremely Divine participation, in a higher
degree than things which merely exist, or which lead an irrational life, or
which are rational like ourselves,'' and so ``at first hand, and under many
forms, participate in the Divine, and, at first hand, and under many forms,
make known the supremely Divine Hiddenness'' (\emph{CH} 4.2). Gregory the
Great places man in the same scale and sees the whole of it gathered in him:
\latin{Habet namque commune esse cum lapidibus, vivere cum arboribus,
sentire cum animalibus, intelligere cum angelis. Si ergo commune habet
aliquid cum omni creatura homo, juxta aliquid omnis creatura est homo} ---
man has being in common with stones, life with trees, sensation with the
animals, understanding with the angels; and since he has something in
common with every creature, in some way every creature is man (\emph{Hom.\
in Evang.} 29.2). Augustine names man's place in the order and its promise:
God created man, ``whose nature was to be a mean between the angelic and
bestial,'' so that if he kept God's commandments ``he should pass into the
company of the angels, and obtain, without the intervention of death, a
blessed and endless immortality'' (\emph{De civitate Dei} XII.21). The
angels are the summit toward which the visible creation leans, and man, who
understands with the angels, is the creature in whom the lower world is
carried up to them.

\subsection{The Channel of the Divine Light}

The order that makes the world one is an order of causes, and the angels
are the highest causes God has made. The \emph{Contra gentiles} gives the
reason in a sentence that could stand over the treatise on government:
\latin{Ad dignitatem regentis pertinet ut habeat multos ministros, et
diversos sui regiminis executores: quia tanto altius et maius ostendetur
suum dominium, quanto plures in diversis gradibus ei subduntur} --- it
belongs to the dignity of a ruler to have many ministers and diverse
executors of his rule, since his dominion is shown the higher and greater the
more there are who are subject to him in diverse grades; and no dignity of any
ruler is comparable with the divine (\emph{SCG} III.77). The order of
causes is nobler than the order of effects, as a cause is nobler than its
effect, \latin{si autem non essent aliquae causae mediae exequentes divinam
providentiam non esset in rebus ordo causarum, sed effectuum tantum}
(ibid.). So God, who orders even the least things immediately by his wisdom,
executes his providence through intermediate causes, and rules the lower
creation through the intellectual creatures: \latin{sicut supremae creaturae
sunt sub Deo et gubernantur ab ipso, ita inferiores creaturae sint sub
superioribus et regantur ab ipsis} (\emph{SCG} III.78); and among the
intellectual creatures themselves the lower are governed by the higher
(\emph{SCG} III.79). ``Therefore God so governs things that He makes some of
them to be causes of others in government; as a master, who not only imparts
knowledge to his pupils, but gives also the faculty of teaching others''
(\ST{I q.~103 a.~6}).

Dionysius draws the procession of that light in images that the Church has
never ceased to use. ``Even as our sun---not as calculating or choosing, but
by its very being, enlightens all things able to partake of its light in
their own degree---so too the Good \dots\ by Its very existence sends to all
things that be, the rays of Its whole goodness, according to their
capacity.'' The intelligible essences receive those rays first; they ``are
illuminated as to the reasons of things, in a manner peculiar to themselves;
and they again convey to their kindred spirits things appropriate to them''
(\emph{DN} 4.1). The purpose of every hierarchy is to make its members
``Divine images, mirrors most luminous and without flaw, receptive of the
primal light and the supremely Divine ray, and devoutly filled with the
entrusted radiance, and again, spreading this radiance ungrudgingly to those
after it'' (\emph{CH} 3.2). So the light passes from the first hierarchy,
``and from this (Order) again, in due degree, the second, and from the
second, the third, and from the third, our Hierarchy, is reverently conducted
to the super-original Origin and End of all good order, according to the
self-same law of well-ordered regularity, in Divine harmony and proportion''
(\emph{CH} 10.1). The \emph{Contra gentiles} names the law by which the
light remains one through all its steps: \latin{in vi superioris virtutis
omnes inferiores agunt} (\emph{SCG} III.80), so that what belongs to the
Seraphim every lower order carries out in their power; and the \emph{Summa}
names the reason for the steps themselves, that each higher angel divides
the knowledge it receives universally, as a teacher does, so that the lower
may grasp it (\ST{I q.~106 a.~1}; \S\ref{sec:q106}).

Augustine, meditating on the government of the world in the third book
\emph{On the Trinity}, gives the whole vision in one passage, from which
Aquinas takes the sentence he sets over God's government through creatures,
over the ordering of the orders, and over the rule of the good angels over
the bad (\ST{I q.~103 a.~6 s.c.}; \S\S\ref{sec:q108},~\ref{sec:q109}).
``For there the will of God, who makes His angels spirits, and His ministers
a flaming fire, presiding among spirits which are joined in perfect peace and
friendship, and combined in one will by a kind of spiritual fire of charity,
as it were in an elevated and holy and secret seat, as in its own house and
in its own temple, thence diffuses itself through all things by certain most
perfectly ordered movements of the creature; first spiritual, then corporeal;
and uses all according to the unchangeable pleasure of its own purpose \dots\
For nothing is done visibly or sensibly, unless either by command or
permission from the interior palace, invisible and intelligible, of the
supreme Governor, according to the unspeakable justice of rewards and
punishments, of favor and retribution, in that far-reaching and boundless
commonwealth of the whole creature'' (\emph{De Trinitate} III.4.9). The
angels are that palace's court, the first movements of its will, and the
first citizens of its commonwealth. The angels' government of bodies and of
men, their missions, and their guardianship are treated at
\S\S\ref{sec:government}--\ref{sec:guardians}.

The light passes through the angels without being diminished into them.
Bernard, having named the orders, turns each toward the God it shows: ``The
Seraphim burn, but with the fire of God, or rather with fire for God \dots\
The Dominions rule, but they rule in subjection to a ruler, and serve Him as
well \dots\ The Principalities lead and govern; but they themselves are
governed: so that they would no longer know how to govern, if they ceased to
be governed \dots\ Angels and Archangels are with us, but He is more our own
who is not only with us but in us'' (\emph{De consideratione} V.5.11). God
alone acts from within the will and justifies (\ST{I q.~106 a.~2});
the angels enlighten and persuade. Their whole greatness is to be
transparent: ``Wherefore, beyond all, they are deemed pre-eminently worthy of
the appellation Angelic, on the ground that the supremely Divine
illumination comes to them at first hand, and, through them, there pass to
us manifestations above us'' (\emph{CH} 4.2).

\subsection{Charity, the Bond of the Order}

The hierarchy is an order of love as much as of light. Dionysius teaches that
divine love ``is extatic, not permitting (any) to be lovers of themselves,
but of those beloved. They shew this too, the superior by becoming mindful
of the inferior; and the equals by their mutual coherence; and the inferior,
by a more divine respect towards things superior'' (\emph{DN} 4.13); and he
quotes from the hymns of Hierotheus the definition that sums up the whole
angelic order: ``Love, whether we speak of Divine, or Angelic, or intelligent,
or psychical, or physical, let us regard as a certain unifying and combining
power, moving the superior to forethought for the inferior, and the equals
to a mutual fellowship, and lastly, the inferior to respect towards the
higher and superior'' (\emph{DN} 4.15). The same love is God's own. It is
``a loving Movement, simplex, self-moved, self-operating, pre-existing in the
Good, and from the Good bubbling forth to things existing, and again
returning to the Good, in which also the Divine Love indicates distinctly Its
own unending and unbeginning, as it were a sort of everlasting circle
whirling round in unerring combination, by reason of the Good, from the
Good, and in the Good, and to the Good'' (\emph{DN} 4.14). The circle of the
\emph{Contra gentiles}, by which the universe is perfected in returning to its
principle, is the circle of this love.

Aquinas gives the metaphysical form of the same truth when he asks whether
the higher angels share all they know: ``it is of the nature of good to
communicate itself to others,'' and ``much more therefore do the holy angels,
who enjoy the plenitude of participation of the Divine goodness, impart the
same to those below them'' (\ST{I q.~106 a.~4}). Gregory gives it its
evangelical form: in the highest city certain things are special to each and
yet common to all, \latin{quia per caritatem Spiritus ab alio in aliis
habentur} (\emph{Hom.\ in Evang.} 34.14). Grade in heaven is no barrier to
communion, for ``each gift is more perfectly possessed by the one who can
communicate it'' (\ST{I q.~108 a.~2 ad~2}); the higher are higher in their
power to give. Gregory's image of the jewels in the King's ornament shows
the other side. The stones were pierced in the day of their creation, ``because,
namely, he was created capable of love''; the one who scorned the gold of
charity fell, and the others, ``bound together by charity mutually
penetrating them,'' received it as a gift ``that they should now be never
loosened by falling from the ornament of the King'' (\emph{Moralia}
XXXII.23.48). Even the fallen remain within the order by their unlost
nature and are used by God for his purposes (\ST{I q.~109 aa.~1,~4};
\S\ref{sec:q109}). The peace of the whole is the peace Augustine names: ``in
heaven there is unbroken peace, both between all the intelligent creatures
that exist there, and between these and their Creator'' (\emph{Enchiridion}
63).

\subsection{The Ranks Filled}

The order of the universe is not finished while the angelic ranks await
their companions. Gregory reads it in the parable of the lost drachma: the
woman had ten, \latin{quia novem sunt ordines angelorum. Sed ut compleretur
electorum numerus, homo decimus est creatus} (\emph{Hom.\ in Evang.} 34.6).
The heavenly city is made of angels and men, \latin{ad quam tantum credimus
humanum genus ascendere, quantos illic contigit electos angelos remansisse}
(34.11), and men enter it in the ranks whose lives they have imitated: the
lowly who announce little things lovingly, among the Angels; those who burn
with contemplation and kindle others, among the Seraphim (34.11;
\S\ref{sec:q108}). Aquinas holds with the Master that the elect are taken
into the nine orders themselves, higher and lower according to their merits,
and not into a tenth order of their own; ``by the gift of grace men can merit
glory in such a degree as to be equal to the angels, in each of the angelic
grades'' (\ST{I q.~108 a.~8}), as the Lord promised that the children of the
resurrection ``are equal to the angels and are the children of God'' (Luke
20:36). The Church on earth and the Church in heaven make one hierarchy:
\latin{unde in patria non erit alia hierarchia hominum et Angelorum, sed una
et eadem et homines in ordines Angelorum distribuentur} (\emph{In II Sent.}
d.~9 q.~1 a.~8 ad~4); and the same commentary sets the Blessed Virgin above
them all (\S\ref{sec:christ-mary}).

Augustine sees the filling of the ranks as the fruit of the Cross. ``Now it
was not for the angels that Christ died. Yet what was done for the
redemption of man through His death was in a sense done for the angels,
because the enmity which sin had put between men and the holy angels is
removed, and friendship is restored between them, and by the redemption of
man the gaps which the great apostasy left in the angelic host are filled
up'' (\emph{Enchiridion} 61). ``And, of course, the holy angels, taught by
God, in the eternal contemplation of whose truth their happiness consists,
know how great a number of the human race are to supplement their ranks, and
fill up the full tale of their citizenship. Wherefore the apostle says, that
`all things are gathered together in one in Christ, both which are in heaven
and which are on earth.' The things which are in heaven are gathered together
when what was lost therefrom in the fall of the angels is restored from among
men'' (\emph{Enchiridion} 62). The Apostle's words in the Douay are that God
purposed ``to re-establish all things in Christ, that are in heaven and on
earth, in him'' (Eph 1:10), and ``through him to reconcile all things unto
himself, making peace through the blood of his cross, both as to the things
that are on earth and the things that are in heaven'' (Col 1:20). The number
of the elect Aquinas refers, with the Fathers, to God alone: \latin{ille scit cui
soli cognitus est numerus electorum, in superna felicitate locandus}
(\emph{In II Sent.} d.~9 q.~1 a.~8).

Already the faithful stand within this one city. ``But you are come to mount
Sion and to the city of the living God, the heavenly Jerusalem, and to the
company of many thousands of angels, and to the church of the firstborn who
are written in the heavens, and to God the judge of all, and to the spirits
of the just made perfect'' (Heb 12:22--23). Augustine calls the angels and
the saints together ``one city of God---His living sacrifice, and His living
temple,'' of which a part ``is hereafter to be united to the immortal angels,
and which at present is being gathered from among mortal men'' (\emph{De
civitate Dei} XII.9). ``This part of the Church, then, which is made up of
the holy angels and the hosts of God, shall become known to us in its true
nature, when, at the end of the world, we shall be united with it in the
common possession of everlasting happiness'' (\emph{Enchiridion} 61).

\subsection{The Praise of the Whole Creation}

The order that descends as light returns as praise. At the beginning of the
world, ``the morning stars praised me together, and all the sons of God made
a joyful melody'' (Job 38:7). Isaiah saw the Lord ``sitting upon a throne
high and elevated,'' and the Seraphim ``cried one to another, and said: Holy,
holy, holy, the Lord God of hosts, all the earth is full of his glory'' (Isa
6:1--3). Dionysius teaches that the Word of God ``has transmitted its hymns to
those on earth, in which are Divinely shewn the excellency of its most
exalted illumination. For some of its members, to speak after sensible
perception, proclaim as a `voice of many waters,' `Blessed is the glory of
the Lord from His place' and others cry aloud that frequent and most august
hymn of God, `Holy, Holy, Holy, Lord of Sabaoth, the whole earth is full of
His glory''' (\emph{CH} 7.4). The first hierarchy, he says, ``dances round
His eternal knowledge in the most exalted ever-moving stability,'' and
teaches those after it ``That it is just and right that the august
Godhead---Itself both above praise, and all-praiseworthy---should be known and
extolled by the God-receptive minds'' (ibid.).

Aquinas says that ``the angels are ever speaking to God in the sense of
praising and admiring Him and His works'' (\ST{I q.~107 a.~3 ad~2}), and
Gregory that ``the voice of the Angels in the praises of God is the very
admiration itself of inward contemplation'' (\emph{Moralia} II.10;
\S\ref{sec:q107}). This praise does not end when the world ends. The orders
remain for ever by nature, grace, and glory, and their offices, which in
this age lead others to the end, remain as they befit the end attained:
``Thus also the various ranks of soldiers have different duties to perform in
battle and in triumph'' (\ST{I q.~108 a.~7}). John saw the triumph. Round the throne the
living creatures ``rested not day and night, saying: Holy, Holy, Holy, Lord
God Almighty, who was and who is and who is to come,'' and the ancients cast
their crowns before the throne: ``Thou art worthy, O Lord our God, to
receive glory and honour and power. Because thou hast created all things:
and for thy will they were and have been created'' (Apoc 4:8--11). ``And I
beheld, and I heard the voice of many angels round about the throne and the
living creatures and the ancients (and the number of them was thousands of
thousands), Saying with a loud voice: The Lamb that was slain is worthy to
receive power and divinity and wisdom and strength and honour and glory and
benediction. And every creature which is in heaven and on the earth and
under the earth, and such as are in the sea, and all that are in them, I
heard all saying: To him that sitteth on the throne and to the Lamb,
benediction and honour and glory and power, for ever and ever'' (Apoc
5:11--13). The angels begin the song and the whole creation takes it up,
through the angels, through man in whom every creature is gathered, through
the Church which on earth joins the Seraphim's hymn in her liturgy
(\S\ref{sec:liturgy}). The Psalmist calls the whole order to it at once:
``Bless the Lord, all ye his angels: you that are mighty in strength, and
execute his word, hearkening to the voice of his orders. Bless the Lord, all
ye his hosts: you ministers of his that do his will. Bless the Lord, all his
works: in every place of his dominion, O my soul, bless thou the Lord'' (Ps
102[103]:20--22).

This is the end toward which the whole treatise has moved. The universe
proceeds from God's goodness in many unequal grades; its highest grade is the
intellectual creature, which alone can return to God by knowing and loving
him; the angels, first in that return, receive the light first and pass it
on, and bind the whole in charity; men are taken up into their ranks; and at
the last, when the Son has delivered up the kingdom to the Father, ``God may
be all in all'' (1 Cor 15:28). Gregory, having led his people through the
nine orders of the heavenly citizens, turned them back to their own lives
with a sentence that closes the contemplation: \latin{Suspiremus ergo ad
eos, de quibus loquimur, sed redeamus ad nos} --- let us sigh for those of
whom we speak, but let us return to ourselves (\emph{Hom.\ in Evang.}
34.15). ``Praise ye him, all his angels, praise ye him, all his hosts''
(Ps 148:2).
